\documentclass[epsfig,12pt]{article}
\usepackage{epsfig}
\usepackage{graphicx}
\usepackage{rotating}
\usepackage{latexsym}
\usepackage{amsmath}
\usepackage{amssymb}
\usepackage{relsize}
\usepackage{geometry}
\geometry{letterpaper}
\usepackage{color}
\usepackage{bm}
\usepackage{slashed}
\usepackage[normalem]{ulem}
%\usepackage{showlabels}




%%%%%%%%%%%%%%%%%%%%%%%%%%%%%%%%%%%%%%%%%%%%%%%%%%%%%%%%%%%%%%%%%%%%%%%%%%%%%%%%
%                                                                              %
%                                                                              %
%                     D O C U M E N T   S E T T I N G S                        %
%                                                                              %
%                                                                              %
%%%%%%%%%%%%%%%%%%%%%%%%%%%%%%%%%%%%%%%%%%%%%%%%%%%%%%%%%%%%%%%%%%%%%%%%%%%%%%%%
\def\baselinestretch{1.1}
\renewcommand{\theequation}{\thesection.\arabic{equation}}

\hyphenation{con-fi-ning}
\hyphenation{Cou-lomb}
\hyphenation{Yan-ki-e-lo-wicz}




%%%%%%%%%%%%%%%%%%%%%%%%%%%%%%%%%%%%%%%%%%%%%%%%%%%%%%%%%%%%%%%%%%%%%%%%%%%%%%%%
%                                                                              %
%                                                                              %
%                      C O M M O N   D E F I N I T I O N S                     %
%                                                                              %
%                                                                              %
%%%%%%%%%%%%%%%%%%%%%%%%%%%%%%%%%%%%%%%%%%%%%%%%%%%%%%%%%%%%%%%%%%%%%%%%%%%%%%%%
\def\beq{\begin{equation}}
\def\eeq{\end{equation}}
\def\beqn{\begin{eqnarray}}
\def\eeqn{\end{eqnarray}}
\def\beqn{\begin{eqnarray}}
\def\eeqn{\end{eqnarray}}
\def\nn{\nonumber}
\def\ba{\beq\new\begin{array}{c}}
\def\ea{\end{array}\eeq}
\def\be{\ba}
\def\ee{\ea}


\newcommand{\nfour}{${\cal N}=4\;$}
\newcommand{\none}{${\mathcal N}=1\,$}
\newcommand{\nonen}{${\mathcal N}=1$}
\newcommand{\ntwo}{${\mathcal N}=2$}
\newcommand{\ntt}{${\mathcal N}=(2,2)\,$}
\newcommand{\nzt}{${\mathcal N}=(0,2)\,$}
\newcommand{\ntwon}{${\mathcal N}=2$}
\newcommand{\ntwot}{${\mathcal N}= \left(2,2\right) $ }
\newcommand{\ntwoo}{${\mathcal N}= \left(0,2\right) $ }
\newcommand{\ntwoon}{${\mathcal N}= \left(0,2\right)$}


\newcommand{\ca}{{\mathcal A}}
\newcommand{\cell}{{\mathcal L}}
\newcommand{\cw}{{\mathcal W}}
\newcommand{\cs}{{\mathcal S}}
\newcommand{\vp}{\varphi}
\newcommand{\pt}{\partial}
\newcommand{\ve}{\varepsilon}
\newcommand{\gs}{g^{2}}
\newcommand{\zn}{$Z_N$}
\newcommand{\cd}{${\mathcal D}$}
\newcommand{\cde}{{\mathcal D}}
\newcommand{\cf}{${\mathcal F}$}
\newcommand{\cfe}{{\mathcal F}}
\newcommand{\ff}{\mc{F}}
\newcommand{\bff}{\ov{\mc{F}}}


\newcommand{\p}{\partial}
\newcommand{\wt}{\widetilde}
\newcommand{\ov}{\overline}
\newcommand{\mc}[1]{\mathcal{#1}}
\newcommand{\md}{\mathcal{D}}
\newcommand{\ml}{\mathcal{L}}
\newcommand{\mw}{\mathcal{W}}
\newcommand{\ma}{\mathcal{A}}


\newcommand{\GeV}{{\rm GeV}}
\newcommand{\eV}{{\rm eV}}
\newcommand{\Heff}{{\mathcal{H}_{\rm eff}}}
\newcommand{\Leff}{{\mathcal{L}_{\rm eff}}}
\newcommand{\el}{{\rm EM}}
\newcommand{\uflavor}{\mathbf{1}_{\rm flavor}}
\newcommand{\lgr}{\left\lgroup}
\newcommand{\rgr}{\right\rgroup}


\newcommand{\Mpl}{M_{\rm Pl}}
\newcommand{\suc}{{{\rm SU}_{\rm C}(3)}}
\newcommand{\sul}{{{\rm SU}_{\rm L}(2)}}
\newcommand{\sutw}{{\rm SU}(2)}
\newcommand{\suth}{{\rm SU}(3)}
\newcommand{\ue}{{\rm U}(1)}


\newcommand{\LN}{\Lambda_\text{SU($N$)}}
\newcommand{\sunu}{{\rm SU($N$) $\times$ U(1) }}
\newcommand{\sunun}{{\rm SU($N$) $\times$ U(1)}}
\def\cfl {$\text{SU($N$)}_{\rm C+F}$ }
\def\cfln {$\text{SU($N$)}_{\rm C+F}$}
\newcommand{\mUp}{m_{\rm U(1)}^{+}}
\newcommand{\mUm}{m_{\rm U(1)}^{-}}
\newcommand{\mNp}{m_\text{SU($N$)}^{+}}
\newcommand{\mNm}{m_\text{SU($N$)}^{-}}
\newcommand{\AU}{\mc{A}^{\rm U(1)}}
\newcommand{\AN}{\mc{A}^\text{SU($N$)}}
\newcommand{\aU}{a^{\rm U(1)}}
\newcommand{\aN}{a^\text{SU($N$)}}
\newcommand{\baU}{\ov{a}{}^{\rm U(1)}}
\newcommand{\baN}{\ov{a}{}^\text{SU($N$)}}
\newcommand{\lU}{\lambda^{\rm U(1)}}
\newcommand{\lN}{\lambda^\text{SU($N$)}}
\newcommand{\bxir}{\ov{\xi}{}_R}
\newcommand{\bxil}{\ov{\xi}{}_L}
\newcommand{\xir}{\xi_R}
\newcommand{\xil}{\xi_L}
\newcommand{\bzl}{\ov{\zeta}{}_L}
\newcommand{\bzr}{\ov{\zeta}{}_R}
\newcommand{\zr}{\zeta_R}
\newcommand{\zl}{\zeta_L}
\newcommand{\nbar}{\ov{n}}
\newcommand{\nnbar}{n\ov{n}}
\newcommand{\muU}{\mu_\text{U}}


\newcommand{\cpn}{CP$^{N-1}$\,}
\newcommand{\CPC}{CP($N-1$)$\times$C }
\newcommand{\CPCn}{CP($N-1$)$\times$C}


\newcommand{\lar}{\lambda_R}
\newcommand{\lal}{\lambda_L}
\newcommand{\larl}{\lambda_{R,L}}
\newcommand{\lalr}{\lambda_{L,R}}
\newcommand{\blar}{\ov{\lambda}{}_R}
\newcommand{\blal}{\ov{\lambda}{}_L}
\newcommand{\blarl}{\ov{\lambda}{}_{R,L}}
\newcommand{\blalr}{\ov{\lambda}{}_{L,R}}


\newcommand{\bgamma}{\ov{\gamma}}
\newcommand{\bpsi}{\ov{\psi}{}}
\newcommand{\bphi}{\ov{\phi}{}}
\newcommand{\bxi}{\ov{\xi}{}}


\newcommand{\qt}{\wt{q}}
\newcommand{\bq}{\ov{q}}
\newcommand{\bqt}{\overline{\widetilde{q}}}


\newcommand{\eer}{\epsilon_R}
\newcommand{\eel}{\epsilon_L}
\newcommand{\eerl}{\epsilon_{R,L}}
\newcommand{\eelr}{\epsilon_{L,R}}
\newcommand{\beer}{\ov{\epsilon}{}_R}
\newcommand{\beel}{\ov{\epsilon}{}_L}
\newcommand{\beerl}{\ov{\epsilon}{}_{R,L}}
\newcommand{\beelr}{\ov{\epsilon}{}_{L,R}}


\newcommand{\bi}{{\bar \imath}}
\newcommand{\bj}{{\bar \jmath}}
\newcommand{\bk}{{\bar k}}
\newcommand{\bl}{{\bar l}}
\newcommand{\bmm}{{\bar m}}


\newcommand{\nz}{{n^{(0)}}}
\newcommand{\no}{{n^{(1)}}}
\newcommand{\bnz}{{\ov{n}{}^{(0)}}}
\newcommand{\bno}{{\ov{n}{}^{(1)}}}
\newcommand{\Dz}{{D^{(0)}}}
\newcommand{\Do}{{D^{(1)}}}
\newcommand{\bDz}{{\ov{D}{}^{(0)}}}
\newcommand{\bDo}{{\ov{D}{}^{(1)}}}
\newcommand{\sigz}{{\sigma^{(0)}}}
\newcommand{\sigo}{{\sigma^{(1)}}}
\newcommand{\bsigz}{{\ov{\sigma}{}^{(0)}}}
\newcommand{\bsigo}{{\ov{\sigma}{}^{(1)}}}


\newcommand{\rrenz}{{r_\text{ren}^{(0)}}}
\newcommand{\bren}{{\beta_\text{ren}}}


\newcommand{\Tr}{\text{Tr}}
\newcommand{\Ts}{\text{Ts}}
\newcommand{\dm}{\hat{{\scriptstyle \Delta} m}}
\newcommand{\dmdag}{\hat{{\scriptstyle \Delta} m}{}^\dag}
\newcommand{\mhat}{\widehat{m}}
\newcommand{\deltam}{{\scriptstyle \Delta} m}
\newcommand{\nvac}{\vec{n}{}_\text{vac}}


\newcommand{\ie}{{\it i.e.}~}
\newcommand{\eg}{{\it e.g.}~}
\newcommand{\ansatz}{{\it ansatz} }


\newcommand{\ii}{\hat\imath}
\newcommand{\jj}{\hat\jmath}
\newcommand{\kk}{\hat k}

\newcommand{\cR}{\mathcal{R}}
\newcommand{\cC}{\mathcal{C}}
\newcommand{\cH}{\mathcal{H}}
\newcommand{\cO}{\mathcal{O}}


% Make the bullet list symbol a small dot
\renewcommand{\labelitemi}{{\Large$\cdot$}}


\begin{document}




%%%%%%%%%%%%%%%%%%%%%%%%%%%%%%%%%%%%%%%%%%%%%%%%%%%%%%%%%%%%%%%%%%%%%%%%%%%%%%%%
%                                                                              %
%                                                                              %
%                            T I T L E   P A G E                               %
%                                                                              %
%                                                                              %
%%%%%%%%%%%%%%%%%%%%%%%%%%%%%%%%%%%%%%%%%%%%%%%%%%%%%%%%%%%%%%%%%%%%%%%%%%%%%%%%
\begin{titlepage}
\ \ 
\vskip 2cm
\begin{center}
{  \Large \bf  Quaternionic Logarithm }
\end{center}
\vskip 0.5cm

\begin{center}

 {\large
 \bf   Pavel A.~Bolokhov
 }
\end {center}

\begin{center}
{\it Theoretical Physics Department, St.Petersburg State University, Ulyanovskaya~1, 
	Peterhof, St.Petersburg, 198504, Russia}
\end{center}




\begin{center}
{\large\bf Abstract}
\end{center}

\hspace{0.3cm}
	We examine the analytical structure of the quaternionic logarithm and
	introduce the concepts of linear integration in quaternionic space, coherent
	with that of the regular complex analysis.
	Along the way, we bring about the suitable definitions of quaternionic derivatives,
	leaving aside the longstanding problem of quaternionic analyticity.
	We briefly discuss possible implications of this analysis for physical scenarios

	We also set the goal of analysis various types of singularities encountered in
	physics from the quaternionic perspective
\vspace{2cm}


\end{titlepage}




%%%%%%%%%%%%%%%%%%%%%%%%%%%%%%%%%%%%%%%%%%%%%%%%%%%%%%%%%%%%%%%%%%%%%%%%%%%%%%%%
%                                                                              %
%                                                                              %
%                            I N T R O D U C T I O N                           %
%                                                                              %
%                                                                              %
%%%%%%%%%%%%%%%%%%%%%%%%%%%%%%%%%%%%%%%%%%%%%%%%%%%%%%%%%%%%%%%%%%%%%%%%%%%%%%%%
\section{Gateway}
\setcounter{equation}{0}

	A typical lecture on introductory mathematical physics often includes 
	a pedagogical derivation of the residue of a holomorphic function,
\[
	\oint\, f(z)\,dz	~~=~~	2\,\pi\,i\,\cdot\,\text{Res}\,f(0)\,.
\]
	This calculation would typically be done in an explicit approach, by
	expanding $ f(z) $ into a series (assuming it is expandable),
\[
	f(z)	~~=~~	\ldots  ~~+~~  \frac{a_{-2}}{z^2}  ~~+~~  \frac{a_{-1}}{z}  ~~+~~  a_0  ~~+~~ a_1\,z  ~~+~~  a_2\,z^2  ~~+~~  a_3\,z^3  ~~+~~  \ldots\,,
\]
	choosing a circular contour around the origin and parametrising it by an 
	angular variable $ \varphi $.
	This way,
\begin{align*}
%
	z^n	& ~~=~~	e^{i\, n\, \varphi}\,,	\\[2mm]
%
	dz	& ~~=~~	i\,z(\varphi)\,d\varphi\,.
\end{align*}
	Then it turns out that the integral of any power of $ z $ vanishes, as $ z $
	is periodic in $ \varphi $,
\[
	i\, a_n\, \int_0^{2\,\pi}\, z^{n \,+\, 1}(\varphi)\, d\varphi	~~=~~	\frac{a_n}{n ~+~ 1}\, z^{n \,+\, 1}(\varphi)\,\bigg|_0^{2\,\pi}	~~=~~	0\,.
\]
	That is --- except when $ n ~=~ -1 $ --- in which case the integrand is a constant.
	This constant then integrates to the celebrated value of $ 2\, \pi\, i $:
\[
	a_{-1}\, \oint\, \frac{dz}{z}	~~=~~	i\, a_{-1}\, \int_0^{2\,\pi}\, d\varphi	~~=~~	2\,\pi\,i \,\cdot\, a_{-1}\,.
\]
	What really happens, and what is easier to explain --- but what is not typically done during such an introductory lecture
	because of the lack of the mathematical basis, which is only being introduced --- is the following.
	Let us take a na\"{i}ve integral of any integer power of $ z $, the same way as we would compute one of a real function of $ x $,
\[
	\int\, z^n\,dz	~~=~~	\frac{1}{n \,+\, 1}\, z^{n \,+\, 1}\,.
\]
	Now, when we evaluate $ z $ at the beginning and at the end of the contour, we realise that we have to ``calculate''
	$ z^{n \,+\, 1} $ at one at the same point (say, $ z ~=~ 1 $).
	And so the integral vanishes.
	But when we pick the power $ 1/z $, this formula no longer works --- the integral no longer yields an integer power of $ z $.
	Instead, it gives us the logarithm,
\[
	\int\, \frac{dz}{z}	~~=~~	\log\,z\,.
\]
	At this point we realise, that the logarithm is \emph{not} evaluated at one and the same point.
	It is instead evaluated at the same position on two different branches of the Riemann surface of the logarithm.
	Hence the $ 2\,\pi\, i $.

	We now re--phrase this conclusion, to emphasize it --- the reason that the integral of $ dz / z $ around a closed contour,
	has given a non--zero result (in contrast to all other powers of $ z $), is that the contour in fact is \emph{not closed}

	A na\"ive student, beginning to study Complex Analysis, would think, that an integral of a (smooth) function should not depend on the \emph{path} of the integration.
	Otherwise, it is not quite clear how to compute integrals.
	And therefore an integral of any function around a closed contour should vanish.
	And we can see that such na\"ive thinking, applied to holomorphic functions, almost works, except for the case of a prime pole, when the contour, strictly speaking,
	is not closed


%%%%%%%%%%%%%%%%%%%%%%%%%%%%%%%%%%%%%%%%%%%%%%%%%%%%%%%%%%%%%%%%%%%%%%%%%%%%%%%%
%%%%%%%%%%%%%%%%%%%%%%%%%%%%%%%%%%%%%%%%%%%%%%%%%%%%%%%%%%%%%%%%%%%%%%%%%%%%%%%%
\vskip 0.8cm
\centerline{*\qquad\qquad\qquad*\qquad\qquad\qquad*}
\vskip 0.6cm
%%%%%%%%%%%%%%%%%%%%%%%%%%%%%%%%%%%%%%%%%%%%%%%%%%%%%%%%%%%%%%%%%%%%%%%%%%%%%%%%
%%%%%%%%%%%%%%%%%%%%%%%%%%%%%%%%%%%%%%%%%%%%%%%%%%%%%%%%%%%%%%%%%%%%%%%%%%%%%%%%


	Of course, everything that we have just described has been known for almost two hundred years.
	However, one can still ask some quite na\"ive questions.
	We can state, that the integral of any holomorphic function around a closed contour vanishes,
	provided that the contour is in fact closed.
	The reason that functions have residues, is that when we try to encircle their poles by a contour,
	we end up with a non--closed contour.
	The contour
\[
\]
	around the origin does not end at the same point where it began (we will talk about why, in a moment).
\[
	\text{But is it possible to \emph{close} this contour, such that the integral along it \emph{would} vanish?}
\]
	Now, indeed, why do we say that the contour around the origin is not closed?
	To this end, function $ f(z) ~=~ 1 / z $ is smooth and single--valued, it does \emph{not} have any branch cuts.
	Yet, when we calculate its integral we end up with a multi--branch function.
	How does that happen?

	Well, there is nothing mysterious.
	When we say that we ``calculate'' the integral
\[	
	\int\, z^n\, dz\,,
\]
	we in fact imply that the 1--form $ z^{n}\,dz $ is exact, and,
\[
	z^n\,dz	~~=~~	\frac{1}{n \,+\, 1}\, d\,z^{n \,+\, 1}\,.
\]
	Then, applying integration to both sides, we end up with the function $ z^{n \,+\, 1} $.
	This argument does not pass through with the function $ 1 / z $.
	The 1--form $ dz/z $ is \emph{closed}, but not \emph{exact},
\[
	\frac{dz}{z}	~~\neq~~	d\,\log\,z\,.
\]
	Locally, these two forms are equal everywhere (except for the branch cut), but not globally.
	They cannot possibly be equal globally, just for the fact that the logarithm is a multi--valued function,
	and has a branch cut (which, although not a singularity of course, in this comparison can be viewed as a ``discontinuity'').
	Function $ dz / z $, in contrast, is perfectly smooth and defined everywhere on the normal complex plane, 
	just less the origin.

	This means, that the form $ dz / z $ in the equality
\beq
\label{int-dz-z}
	\int\, \frac{dz} z	~~=~~	\log\,z\,,
\eeq
	is not entirely the same as one would normally think,
	but rather, its extension to the Riemann surface of the logarithm.
	With this clarification, the equality \eqref{int-dz-z} is correct.
	So, we have had to \emph{extend} function $ 1 / z $ to the Riemann surface of $ \log z $,
	and only then could we perform the integration.

	The following questions arise, 
\begin{align*}
	& \text{did function $ 1 / z $ in fact know that its integral would be multi--valued?}\\[2mm]
	& \text{did it have that information encoded anywhere?}
\end{align*}

	It turns out that yes, in a way, it did know that.
	As for question (??) --- it \emph{is} possible to close the contour $ \gamma $ in such a way,
	that the integral around it would vanish.
	In order to do that, however, we find that we need to step out of the complex plane.
	On that path, we realise that we are entering a whole new world.


\newpage
%%%%%%%%%%%%%%%%%%%%%%%%%%%%%%%%%%%%%%%%%%%%%%%%%%%%%%%%%%%%%%%%%%%%%%%%%%%%%%%%
%                                                                              %
%                                                                              %
%                             A N A L Y T I C I T Y                            %
%                                                                              %
%                                                                              %
%%%%%%%%%%%%%%%%%%%%%%%%%%%%%%%%%%%%%%%%%%%%%%%%%%%%%%%%%%%%%%%%%%%%%%%%%%%%%%%%
\section{Analyticity}
\label{section-dirac}
\setcounter{equation}{0}

	In this section we explain why we do \emph{not} address the issue of quaternionic analiticity in this work.
	The issue becomes apparent during the variable counting

	Consider arbitrary complex functions $ f(z,~ \ov z) $.
	These consist of two real functions of two real variables.
	In terms of the real--valued analysis, a derivative of a complex function is a $ 2 \times 2 $ symmetric real matrix,
\[
	\p f(z, ~\ov z) 	~~=~~	\p \bigl( u ~+~ i\, v \bigr)	~~=~~
		\lgr \begin{matrix}\,\,
			\p_x\, u	&	\p_y\, u	\\[2mm]
			\p_x\, v	&	\p_y\, v
		\,\,\end{matrix} \rgr\!\!.
\]
	This derivative can also as well be described in terms of two complex derivatives --- $ \p_z f $ and $ \p_{\ov z} f $

	The complex analyticity condition
\beq
\label{complex-analyticity}
	\partial_{\ov z}\, f(z,~ \ov z)	~~=~~	 0\,,
\eeq
	imposes two real conditions, leaving us with two independent real functions, which form a single, holomorphic function.
	It makes sense to denote such a function as $ f(z) $ because condition \eqref{complex-analyticity} implies that
	this function is independent of $ \ov z $ as of a linear combination $ x \,-\, i y $.
	This condition also implies that integration of $ f(z) $ can be written and performed in the conventional real--analysis fashion,
\[
	f(b)  ~-~  f(a)		~~=~~	\int_a^b\, \p f			~~=~~
		\int_a^b\, \p_z\,f(z)\, dz\,.
\]

	In contrast, a general quaternionic function consists of four functions of four real variables.
	Therefore, a derivative of such a function is a $ 4 \times 4 $ symmetric matrix
\[
	\p\, f(q)	~~=~~
		\lgr \begin{matrix}\,\,
			\p_t\, f_0	&	\p_x\, f_0	&	...	&	\p_z\,f_1	\\[2mm]
					&	...		&	...	&			\\[2mm]
					&	...		&	...	&			\\[2mm]
			\p_t\,f_3	&	...		&	...	&	\p_z\,f_3
		\,\,\end{matrix} \rgr\!\!,
\]
	which would require \emph{four} quaternionic derivative functions for an equivalent description.
	And this is what we will have to do --- we will introduce four linearly independent quaternionic derivatives.

	But the real problem now is that, unlike in the complex case, one cannot require ``analytical'' functions 
	to be independent from any linear combination other than 
	$ t \,+\, \ii x \,+\, \jj y \,+\, \kk z $.
	This is because variables $ t $, $ x $, $ y $ and $ z $ can each be represented as an ``analytical function'' of $ q $.
	And therefore, any smooth four--component function of real variables $ t $, $ x $, $ y $ and $ z $ must be considered
	an analytical quaternionic function.
	Otherwise the analyticity conditions become too restrictive.

	This does not preclude certain choices of ``preferred'' conditions of analyticity, as we will discuss below,
	but they are still too restrictive in our opinion.

	The problem now is that it is almost meaningless to perform integration of a quaternionic function along a path by itself.
	While this would lead to a certain numerical or algebraic result, it would not be related to what we wish to call an integral of a function.
	In other words,
\[
	\p\,f(q)\,dq	~~\neq~~	df(q)\,.
\]
	Instead, all four linearly independent quaternionic derivatives would have to be involved in order to find the ``change'' of the antideritavive
	as the result,
\[
	df(q)		~~=~~
		\frac{  \p\,f(q)\,dq	~+~	\p_{(\ii)}\,f(q)\,dq^{(\ii)}	~+~	\p_{(\jj)}\,f(q)\,dq^{(\jj)}	~+~	\p_{(\kk)}\,f(q)\,dq^{(\kk)} }
			{  4  }~.
\]
	We call such derivatives \emph{sister} derivatives to the original derivative $ \p f(q) $, and introduce them in the Section~\ref{section-conjugations}.

	In other words, all four sister derivatives need to be known in order to perform quaternionic integration.
	Or, in general, one needs four functions in order to perform integration along a contour.


%%%%%%%%%%%%%%%%%%%%%%%%%%%%%%%%%%%%%%%%%%%%%%%%%%%%%%%%%%%%%%%%%%%%%%%%%%%%%%%%
%                                                                              %
%                          P U R E   F U N C T I O N S                         %
%                                                                              %
%%%%%%%%%%%%%%%%%%%%%%%%%%%%%%%%%%%%%%%%%%%%%%%%%%%%%%%%%%%%%%%%%%%%%%%%%%%%%%%%
\subsection{Pure functions}

	How does this help us with integrating our function $ 1 / z $?
	This function is a derivative of $ \log\, z $, but do we know its sister derivatives?

	Before addressing these questions let us introduce a concept of \emph{pure} functions.
	Any algebraic complex analytical function $ f(z) $ can be easily promoted to a quaternionic function $ f(q) $:
\[
	f(z)	~~=~~	e^z		\qquad\qquad\longrightarrow\qquad\qquad		f(q)	~~=~~	e^q\,.
\]
	We call such functions \emph{pure}, because they do not involve fixed quaternionic constants.
	A more precise definition can be done, but is not necessary for what follows
	It is obvious that any analytical complex function can be promoted this way,
\[
	\sin z	\qquad\longrightarrow\qquad	\sin q\,,\qquad\qquad
	\log z \qquad\longrightarrow\qquad	\log q\,.
\]

	While such functions not necessarily are useful for the quaternionic analysis that we are going to dive into,
	they are useful in relating the quaternionic properties that we find for them to their complex counterparts


	TODO: discuss the mapping of $ i $

\newpage
%%%%%%%%%%%%%%%%%%%%%%%%%%%%%%%%%%%%%%%%%%%%%%%%%%%%%%%%%%%%%%%%%%%%%%%%%%%%%%%%
%                                                                              %
%                                                                              %
%                            C O N J U G A T I O N S                           %
%                                                                              %
%                                                                              %
%%%%%%%%%%%%%%%%%%%%%%%%%%%%%%%%%%%%%%%%%%%%%%%%%%%%%%%%%%%%%%%%%%%%%%%%%%%%%%%%
\section{Conjugations}
\label{section-conjugations}
\setcounter{equation}{0}

	We need to elaborate on the concept of the quaternionic conjugation and introduce the relevant notations.
	Given a quaternion
\[
	q	~~=~~	t  ~~+~~  \ii\, x  ~~+~~  \jj\, y  ~~+~~  \kk\, z\,,
\]
	we denote the usual quaternionic conjugate of $ q $ as
\beq
\label{q-conj}
	\wt q	~~=~~	t  ~~-~~  \ii\, x  ~~-~~  \jj\, y  ~~-~~  \kk\, z\,.
\eeq

	We then introduce directional quaternionic conjugates,
\begin{align}
%
\notag
	q^{(\ii)}	& ~~=~~	-\,\ii\, q \, \ii	~~=~~	t  ~~+~~  \ii\, x  ~~-~~  \jj\, y  ~~-~~  \kk\, z\,,
	\\[2mm]
%
\label{q-dir-conj}
	q^{(\jj)}	& ~~=~~	-\,\ii\, q \, \ii	~~=~~	t  ~~-~~  \ii\, x  ~~+~~  \jj\, y  ~~-~~  \kk\, z\,,
	\\[2mm]
%
\notag
	q^{(\kk)}	& ~~=~~	-\,\ii\, q \, \ii	~~=~~	t  ~~-~~  \ii\, x  ~~-~~  \jj\, y  ~~+~~  \kk\, z\,,
\end{align}
	and in general, for any unit vector $ \hat a $,
\beq
	q^{(\hat a)}	~~=~~	-\, \hat a\, q \, \hat a\,.
\eeq
	The composition of the conjugations \eqref{q-dir-conj} can be readily found,
\[
	\big(\, q^{(\ii)} \,\big)^{(\jj)}	~~=~~	\big(\, q^{(\jj)} \,\big)^{(\ii)}	~~=~~	q^{(\kk)}\,,
	\qquad\qquad
	\big(\, q^{(\ii)} \,\big)^{(\ii)}	~~=~~	q\,.
\]
	These are valid of course as long as $ \ii $, $ \jj $ and $ \kk $ are all orthogonal
	
	We notice that the conjugations \eqref{q-dir-conj} always flip \emph{two} signs.
	In contrast, any composition with the regular quaternionic conjugation \eqref{q-conj} will flip an odd number of signs,
\[
	\overset{\wt{~~~}}{\big(\, q^{(\ii)} \,\big)}		~~=~~	t  ~~-~~  \ii\, x  ~~+~~  \jj\, y  ~~+~~  \kk\, z\,.
\]

	We observe that the conjugations \eqref{q-dir-conj} and \eqref{q-conj} are all completely interchangeable,
\[
	\overset{\wt{~~~}}{\big(\, q^{(\ii)} \,\big)}	~~=~~	\big(\, \wt{q} \,\big)^{(\ii)}\,.
\]
	The following obvious identities can be easily inferred,
\begin{align*}
%
	\big(\, a\, b \,\big)^{(\ii)}	& ~~=~~		a^{(\ii)}\, b^{(\ii)}\,,
	\\[2mm]
%
	\big(\, a\, \wt b \,)^{(\ii)}	& ~~=~~		\wt b{}^{(\ii)}\, \wt a{}^{(\ii)}\,,
	\\[2mm]
%
	a^{(\ii)}\, b~~~		& ~~=~~		\big(\, a\, b^{(\ii)} \big)^{(\ii)}\,.
\end{align*}
	We can also use the directed conjugations to isolate any component of $ q $,
\begin{align}
%
\notag
	t	& ~~=~~	\frac 1 4\, \Big\{\,  q  ~+~  q^{(\ii)}  ~+~  q^{(\jj)}  ~+~  q^{(\kk)} \Big\}\,,
	\\[2mm]
%
\notag
	x	& ~~=~~	\frac{1}{4\,\ii}\, \Big\{\,  q  ~+~  q^{(\ii)}  ~-~  q^{(\jj)}  ~-~  q^{(\kk)} \Big\}\,,
	\\[2mm]
%
\label{comp-from-conj}
	y	& ~~=~~	\frac{1}{4\,\jj}\, \Big\{\,  q  ~-~  q^{(\ii)}  ~+~  q^{(\jj)}  ~-~  q^{(\kk)} \Big\}\,,
	\\[2mm]
%
\notag
	z	& ~~=~~	\frac{1}{4\,\kk}\, \Big\{\,  q  ~-~  q^{(\ii)}  ~-~  q^{(\jj)}  ~+~  q^{(\kk)} \Big\}\,,
\end{align}
	or even express $ \wt{q} $ in terms of them,
\[
	\wt q 	~~=~~ 	\frac 1 4\, \Big\{\,  q  ~-~  q^{(\ii)}  ~-~  q^{(\jj)}  ~-~  q^{(\kk)}  \,\}\,.
\]	
	While being very trivial, identities \eqref{q-dir-conj} and \eqref{comp-from-conj} make it especially obvious,
	that any real--analytical function of coordinates $ t $, $ x $, $ y $ and $ z $
	is essentially an analytical function of quaternions,
	and it is not possible to demand a function to only depend on a certain linear combination of these coordinates
	(which is what is done in the definition of complex analyticity)

	The following quite trivial identity will be found useful,
\[
%
	\text{Re}~ q	~~=~~	\text{Re}~ \wt q	~~=~~	
	\text{Re}~ q^{(\ii)}	~~=~~	\text{Re}~ q^{(\jj)}	~~=~~	\text{Re}~ q^{(\kk)}\,.
\]


%%%%%%%%%%%%%%%%%%%%%%%%%%%%%%%%%%%%%%%%%%%%%%%%%%%%%%%%%%%%%%%%%%%%%%%%%%%%%%%%
%                                                                              %
%                          S C A L A R   P R O D U C T                         %
%                                                                              %
%%%%%%%%%%%%%%%%%%%%%%%%%%%%%%%%%%%%%%%%%%%%%%%%%%%%%%%%%%%%%%%%%%%%%%%%%%%%%%%%
\subsection{Scalar Product}
	
	For what follows, it is necessary to introduce the notations for the ``scalar product'' for quaternions.
	We will be denoting the scalar product both for quaternions and for $ 3 $--vectors using angle brackets,
\begin{align*}
%
	\langle\, a,\, b \,\rangle		& ~~=~~		a_\mu\, b^\mu	~~=~~	
		a^0\, b^0  ~~+~~  a^1\, b^1  ~~+~~  a^2\, b^2  ~~+~~  a^3\, b^3\,,
	\\[2mm]
%
	\langle\, \vec a,\, \vec b \,\rangle	& ~~=~~  a^1\, b^1  ~~+~~  a^2\, b^2  ~~+~~  a^3\, b^3\,,
\end{align*}
	where we now have used the usual index notations for the components of $ 3 $-- and $ 4 $--vectors

	It is trivial to express the scalar product of two quaterions in terms of their conjugates,
\[
	\langle\, a,\, b \,\rangle
	~~=~~	\frac 1 2\, \Big\{\, \wt a\, b  ~+~  \wt b\, a \,\Big\}
	~~=~~	\frac 1 2\, \Big\{\, a\, \wt b  ~+~  b\, \wt a \,\Big\}\,.
\]
	Essentially, the usual quaternionic conjugation turns a contravariant $ 4 $--vector into a covariant one

	Similarly, and suddenly more interestingly, the directed conjugates can be used,
\begin{align}
%
\notag
	\langle\, a,\, b \,\rangle
	& ~~=~~	\frac 1 4\, \Big\{\, \wt a\, b  ~~+~~  \wt a^{(i)}\, b^{(i)}  ~~+~~  \wt a^{(j)}\, b^{(j)}  ~~+~~  \wt a^{(k)}\, b^{(k)} \,\Big\}
	\\[2mm]
%
	& ~~=~~	\frac 1 4\, \Big\{\, a\, \wt b  ~~+~~  a^{(i)}\, \wt b^{(i)}  ~~+~~  a^{(j)}\, \wt b^{(j)}  ~~+~~  a^{(k)}\, \wt b^{(k)} \,\Big\}
	\,.
\end{align}
	These expressions are invariant under any of the conjugations that we have introduced, as well as
	under replacing $ a $ $ \leftrightarrow $ $ b $

	However, if we now introduce the greek--letter notation for the directed conjugates,
\[
	q^{(\mu)}:\qquad	q^{(0)}		~~\equiv~~	q\,,\qquad
	q^{(1)}		~~\equiv~~		q^{(\ii)}\,,\qquad
	q^{(2)}		~~\equiv~~		q^{(\jj)}\,,\qquad
	q^{(3)}		~~\equiv~~		q^{(\kk)}\,,
\]
	the scalar product can suddenly be written in a compact way,
\beq
\label{scalar-product}
	\langle\, a,\, b \,\rangle
	~~=~~	\frac 1 4\, \wt a_{(\mu)}\, b^{(\mu)}
	~~=~~	\frac 1 4\, a^{(\mu)}\, \wt b_{(\mu)}\,.
\eeq
	Since the expression on the left--hand side is real,
	the expressions on the right--hand side do not depend on the order of $ a $ and $ b $,
	despite that all of them are, in fact, quaternions.

	Expression \eqref{scalar-product} is now, unexpectedly, 
	very similar to the usual expression for the scalar product $ a_\mu\, b^\mu $.
	We chose to lower index $ (\mu) $ on $ \wt b $ just for convience and for this resemblance.
	Despite that $ q^{(\mu)} $ looks like ``components'' of vector $ q $, they are, of course, 
	not its components, nor are they even real.
	
	However, it is because of this resemblance that we will be able to manage quaternionic expressions
	involved into integration and differentiation 

\newpage
%%%%%%%%%%%%%%%%%%%%%%%%%%%%%%%%%%%%%%%%%%%%%%%%%%%%%%%%%%%%%%%%%%%%%%%%%%%%%%%%
%                                                                              %
%                                                                              %
%                            D I F F E R E N T I A L                           %
%                                                                              %
%                                                                              %
%%%%%%%%%%%%%%%%%%%%%%%%%%%%%%%%%%%%%%%%%%%%%%%%%%%%%%%%%%%%%%%%%%%%%%%%%%%%%%%%
\section{Differential}
\label{section-dirac}
\setcounter{equation}{0}


\newpage
%%%%%%%%%%%%%%%%%%%%%%%%%%%%%%%%%%%%%%%%%%%%%%%%%%%%%%%%%%%%%%%%%%%%%%%%%%%%%%%%
%                                                                              %
%                                                                              %
%                             D E R I V A T I V E S                            %
%                                                                              %
%                                                                              %
%%%%%%%%%%%%%%%%%%%%%%%%%%%%%%%%%%%%%%%%%%%%%%%%%%%%%%%%%%%%%%%%%%%%%%%%%%%%%%%%
\section{Derivatives}
\label{section-dirac}
\setcounter{equation}{0}


\newpage
%%%%%%%%%%%%%%%%%%%%%%%%%%%%%%%%%%%%%%%%%%%%%%%%%%%%%%%%%%%%%%%%%%%%%%%%%%%%%%%%
%                                                                              %
%                                                                              %
%                               L O G A R I T H M                              %
%                                                                              %
%                                                                              %
%%%%%%%%%%%%%%%%%%%%%%%%%%%%%%%%%%%%%%%%%%%%%%%%%%%%%%%%%%%%%%%%%%%%%%%%%%%%%%%%
\section{Logarithm}
\label{section-dirac}
\setcounter{equation}{0}


\newpage
%%%%%%%%%%%%%%%%%%%%%%%%%%%%%%%%%%%%%%%%%%%%%%%%%%%%%%%%%%%%%%%%%%%%%%%%%%%%%%%%
%                                                                              %
%                                                                              %
%                             I N T E G R A T I O N                            %
%                                                                              %
%                                                                              %
%%%%%%%%%%%%%%%%%%%%%%%%%%%%%%%%%%%%%%%%%%%%%%%%%%%%%%%%%%%%%%%%%%%%%%%%%%%%%%%%
\section{Integration}
\label{section-dirac}
\setcounter{equation}{0}


%%%%%%%%%%%%%%%%%%%%%%%%%%%%%%%%%%%%%%%%%%%%%%%%%%%%%%%%%%%%%%%%%%%%%%%%%%%%%%%%
%                                                                              %
%                                                                              %
%                             C O N C L U S I O N S                            %
%                                                                              %
%                                                                              %
%%%%%%%%%%%%%%%%%%%%%%%%%%%%%%%%%%%%%%%%%%%%%%%%%%%%%%%%%%%%%%%%%%%%%%%%%%%%%%%%
\section{Conclusions}


%%%%%%%%%%%%%%%%%%%%%%%%%%%%%%%%%%%%%%%%%%%%%%%%%%%%%%%%%%%%%%%%%%%%%%%%%%%%%%%%
%                                                                              %
%                                                                              %
%                          A C K N O W L E D G M E N T S                       %
%                                                                              %
%                                                                              %
%%%%%%%%%%%%%%%%%%%%%%%%%%%%%%%%%%%%%%%%%%%%%%%%%%%%%%%%%%%%%%%%%%%%%%%%%%%%%%%%
\section*{Acknowledgements}




%%%%%%%%%%%%%%%%%%%%%%%%%%%%%%%%%%%%%%%%%%%%%%%%%%%%%%%%%%%%%%%%%%%%%%%%%%%%%%%%
%%%%%%%%%%%%%%%%%%%%%%%%%%%%%%%%%%%%%%%%%%%%%%%%%%%%%%%%%%%%%%%%%%%%%%%%%%%%%%%%
%                                                                              %
%                                                                              %
%                                A P P E N D I X                               %
%                                                                              %
%                                                                              %
%%%%%%%%%%%%%%%%%%%%%%%%%%%%%%%%%%%%%%%%%%%%%%%%%%%%%%%%%%%%%%%%%%%%%%%%%%%%%%%%
%%%%%%%%%%%%%%%%%%%%%%%%%%%%%%%%%%%%%%%%%%%%%%%%%%%%%%%%%%%%%%%%%%%%%%%%%%%%%%%%
\pagebreak
\appendix
\setcounter{equation}{0}


%%%%%%%%%%%%%%%%%%%%%%%%%%%%%%%%%%%%%%%%%%%%%%%%%%%%%%%%%%%%%%%%%%%%%%%%%%%%%%%%
%                                                                              %
%                                                                              %
%                                 S P I N O R S                                %
%                                                                              %
%                                                                              %
%%%%%%%%%%%%%%%%%%%%%%%%%%%%%%%%%%%%%%%%%%%%%%%%%%%%%%%%%%%%%%%%%%%%%%%%%%%%%%%%
\section{Spinors}
\label{section-spinors}



%%%%%%%%%%%%%%%%%%%%%%%%%%%%%%%%%%%%%%%%%%%%%%%%%%%%%%%%%%%%%%%%%%%%%%%%%%%%%%%%
%                                                                              %
%                                                                              %
%                            B I B L I O G R A P H Y                           %
%                                                                              %
%                                                                              %
%%%%%%%%%%%%%%%%%%%%%%%%%%%%%%%%%%%%%%%%%%%%%%%%%%%%%%%%%%%%%%%%%%%%%%%%%%%%%%%%
\small
\begin{thebibliography}{99}
\itemsep -2pt

  %\cite{Adler:1988hb}
  \bibitem{Adler:1988hb} 
  S.~L.~Adler,
  %``Quaternionic quantum mechanics and quantum fields,''
  New York, USA: Oxford Univ. Pr. (1995) 586 p. (International series of monographs on physics, 88)

  %\cite{Maxwell}
  \bibitem{Maxwell}
  J.~C.~Maxwell,
  {\it A Treatise on Electricity and Magnetism},
  Oxford, Clarendon Press (1873)

  %\cite{Tait}
  \bibitem{Tait}
  P.~G.~Tait,
  {\it An Elementary Treatise on Quaternions},
  Cambridge University Press (1890)

  %\cite{Finkelstein:1961tk}
  \bibitem{Finkelstein:1961tk} 
  D.~Finkelstein, J.~M.~Jauch and D.~Speiser,
  %``Foundations of quaternion quantum mechanics,''
  J.\ Math.\ Phys.\  {\bf 3}, 207 (1962).
  doi:10.1063/1.1703794
  %%CITATION = doi:10.1063/1.1703794;%%

  %\cite{Finkelstein:1962tj}
  \bibitem{Finkelstein:1962tj} 
  D.~Finkelstein, J.~M.~Jauch, S.~Schminovich and D.~Speiser,
  %``Principle of general Q covariance,''
  J.\ Math.\ Phys.\  {\bf 4}, 788 (1963).
  doi:10.1063/1.1724320
  %%CITATION = doi:10.1063/1.1724320;%%

  %\cite{Edmonds:1972vd}
  \bibitem{Edmonds:1972vd} 
  J.~D.~Edmonds,
  %``Nature's natural numbers - relativistic quantum theory over the ring of complex quaternions,''
  Int.\ J.\ Theor.\ Phys.\  {\bf 6}, 205 (1972).
  doi:10.1007/BF00672074
  %%CITATION = doi:10.1007/BF00672074;%%

  %\cite{Edmonds:1973pp}
  \bibitem{Edmonds:1973pp} 
  J.~D.~Edmonds,
  %``Possible baryon conservation from quaternion charge,''
  Lett.\ Nuovo Cim.\  {\bf 7S2}, 398 (1973)
  [Lett.\ Nuovo Cim.\  {\bf 7}, 398 (1973)].
  doi:10.1007/BF02735143
  %%CITATION = doi:10.1007/BF02735143;%%

  %\cite{Edmonds:1974yq}
  \bibitem{Edmonds:1974yq} 
  J.~D.~Edmonds,
  %``Quaternion wave equations in curved space-time,''
  Int.\ J.\ Theor.\ Phys.\  {\bf 10}, 115 (1974).
  doi:10.1007/BF01810397
  %%CITATION = doi:10.1007/BF01810397;%%

  %\cite{Gunaydin:1973rs}
  \bibitem{Gunaydin:1973rs} 
  M.~Gunaydin and F.~Gursey,
  %``Quark structure and octonions,''
  J.\ Math.\ Phys.\  {\bf 14}, 1651 (1973).
  doi:10.1063/1.1666240
  %%CITATION = doi:10.1063/1.1666240;%%

  %\cite{Gunaydin:1974fb}
  \bibitem{Gunaydin:1974fb} 
  M.~Gunaydin and F.~Gursey,
  %``Quark Statistics and Octonions,''
  Phys.\ Rev.\ D {\bf 9}, 3387 (1974).
  doi:10.1103/PhysRevD.9.3387
  %%CITATION = doi:10.1103/PhysRevD.9.3387;%%

  %\cite{Gough:1986}
  \bibitem{Gough:1986}
  W.~Gough,
  %``The analysis of spin and spin-orbit coupling in quantum and classical physics by quaternions,''
  Eur.\ J.\ Phys.\ {\bf 7}, 35 (1986).
  doi:10.1088/0143-0807/7/1/007
  %%CITATION = doi:10.1088/0143-0807/7/1/007;%%

  %\cite{Gough:1989}
  \bibitem{Gough:1989}
  W.~Gough,
  %``A quaternion expression for the quantum mechanical probability and current densities,''
  Eur.\ J.\ Phys.\ {\bf 10}, 188 (1989).
  doi:10.1088/0143-0807/10/3/005
  %%CITATION = doi:10.1088/0143-0807/10/3/005;%%

  %\cite{DeLeo:1995ww}
  \bibitem{DeLeo:1995ww} 
  S.~De Leo,
  %``One component Dirac equation,''
  Int.\ J.\ Mod.\ Phys.\ A {\bf 11}, 3973 (1996)
  doi:10.1142/S0217751X96001863
  [hep-th/9508010].
  %%CITATION = doi:10.1142/S0217751X96001863;%%

  %\cite{DeLeo:1995yq}
  \bibitem{DeLeo:1995yq} 
  S.~De Leo and P.~Rotelli,
  %``The Quaternionic Dirac Lagrangian,''
  Mod.\ Phys.\ Lett.\ A {\bf 11}, 357 (1996)
  doi:10.1142/S0217732396000400
  [hep-th/9509059].
  %%CITATION = doi:10.1142/S0217732396000400;%%

  %\cite{DeLeo:1997yj}
  \bibitem{DeLeo:1997yj} 
  S.~De Leo and W.~A.~Rodrigues, Jr.,
  %``Quaternionic electron theory: Geometry, algebra and Dirac's spinors,''
  Int.\ J.\ Theor.\ Phys.\  {\bf 37}, 1707 (1998)
  doi:10.1023/A:1026692508708
  [hep-th/9806058].
  %%CITATION = doi:10.1023/A:1026692508708;%%

  %\cite{DeLeo:1998yi}
  \bibitem{DeLeo:1998yi} 
  S.~De Leo and W.~A.~Rodrigues, Jr.,
  %``Quaternionic electron theory: Dirac's equation,''
  Int.\ J.\ Theor.\ Phys.\  {\bf 37}, 1511 (1998)
  doi:10.1023/A:1026611718277
  [hep-th/9806057].
  %%CITATION = doi:10.1023/A:1026611718277;%%

  %\cite{DeLeo:2000ik}
  \bibitem{DeLeo:2000ik}
  S.~De Leo,
  %``Quaternionic Lorentz group and Dirac equation,''
  Found.\ Phys.\ Lett.\  {\bf 14}, no. 1, 37 (2001)
  doi:10.1023/A:1012077227985
  [hep-th/0103129].
  %%CITATION = doi:10.1023/A:1012077227985;%%

  %\cite{Furey:2010fm}
  \bibitem{Furey:2010fm} 
  C.~Furey,
  %``Unified Theory of Ideals,''
  Phys.\ Rev.\ D {\bf 86}, 025024 (2012)
  doi:10.1103/PhysRevD.86.025024
  [arXiv:1002.1497 [hep-th]].
  %%CITATION = doi:10.1103/PhysRevD.86.025024;%%

  %\cite{Furey:2014iwa}
  \bibitem{Furey:2014iwa} 
  C.~Furey,
  %``Generations: Three Prints, in Colour,''
  JHEP {\bf 1410}, 046 (2014)
  doi:10.1007/JHEP10(2014)046
  [arXiv:1405.4601 [hep-th]].
  %%CITATION = doi:10.1007/JHEP10(2014)046;%%

  %\cite{Furey:2015tqa}
  \bibitem{Furey:2015tqa} 
  C.~Furey,
  %``Charge quantization from a number operator,''
  Phys.\ Lett.\ B {\bf 742}, 195 (2015)
  doi:10.1016/j.physletb.2015.01.023
  [arXiv:1603.04078 [hep-th]].
  %%CITATION = doi:10.1016/j.physletb.2015.01.023;%%

  %\cite{Bisht:2007ju}
  \bibitem{Bisht:2007ju} 
  P.~S.~Bisht and O.~P.~S.~Negi,
  %``Revisiting quaternionic dual electrodynamics,''
  Int.\ J.\ Theor.\ Phys.\  {\bf 47}, 3108 (2008)
  doi:10.1007/s10773-008-9744-8
  [arXiv:0709.0088 [hep-th]].
  %%CITATION = doi:10.1007/s10773-008-9744-8;%%

  %\cite{Rawat:2011ch}
  \bibitem{Rawat:2011ch} 
  A.~S.~Rawat and O.~P.~S.~Negi,
  %``Quaternion Gravi-Electromagnetism,''
  Int.\ J.\ Theor.\ Phys.\  {\bf 51}, 738 (2012)
  doi:10.1007/s10773-011-0953-1
  [arXiv:1107.0916 [physics.gen-ph]].
  %%CITATION = doi:10.1007/s10773-011-0953-1;%%  

  %\cite{Chanyal:2012rj}
  \bibitem{Chanyal:2012rj} 
  B.~C.~Chanyal, P.~S.~Bisht, T.~Li and O.~P.~S.~Negi,
  %``Octonion Quantum Chromodynamics,''
  Int.\ J.\ Theor.\ Phys.\  {\bf 51}, 3410 (2012)
  doi:10.1007/s10773-012-1222-7
  [arXiv:1204.0242 [physics.gen-ph]].
  %%CITATION = doi:10.1007/s10773-012-1222-7;%%

  %\cite{thesis}
  \bibitem{thesis}
  C.~Furey,
  %``Standard model physics from an algebra?,''
  arXiv:1611.09182 [hep-th].
  %%CITATION = ARXIV:1611.09182;%%

  %\cite{girard}
  \bibitem{girard}
  P.~Girard,
  {\it Quaternions, Clifford algebras and relativistic physics},
  Birkh\"auser Verlag (2007).

  %\cite{sudbery}
  \bibitem{sudbery}
  A.~Sudbery,
  {\it Quaternionic Analysis},
  Math.\ Proc.\ Camb.\ Math.\ Soc.\  {\bf 85} (1979), 199-225.

  %\cite{frenkel}
  \bibitem{frenkel}
  I.~Frenkel, M.~Libine,
  {\it Quaternionic analysis, representation theory and physics},
  Adv.\ Math.\  {\bf 218} (2008), 1806-1877
  doi:10.1016/j.aim.2008.03.021
  %%CITATION = doi:10.1016/j.aim.2008.03.021;%%

  %\cite{Berestetsky:1982aq}
  \bibitem{Berestetsky:1982aq} 
  V.~B.~Berestetskii, E.~M.~Lifshitz and L.~P.~Pitaevskii,
  {\it Quantum Electrodynamics},
  2nd ed.,
  Oxford, New York: Pergamon Press, 1982.
  %%CITATION = INSPIRE-185062;%%
\end{thebibliography}


\end{document}
